\documentclass{article}
\usepackage{graphicx} % Required for inserting images
\usepackage{caption} % Required for \captionof
\usepackage{amsmath}
\usepackage{float}
\usepackage{comment}
\usepackage[T1]{fontenc}
\usepackage[utf8]{inputenc}
\usepackage{lmodern}
\usepackage{microtype}
\usepackage{xcolor}
\usepackage{framed}
\renewenvironment{leftbar}{%
  \def\FrameCommand{\textcolor{orange}{\vrule width 3pt} \hspace{10pt}}%
  \MakeFramed {\advance\hsize-\width \FrameRestore}}%
 {\endMakeFramed}

\title{Mechanikai mérés}
\author{Kern Luca, Vincze Csongor}
\date{2026 Április 22}

\begin{document}

\maketitle

%%%%%%%%%%%%%%%%%%%%%%%%%%
% 1. FELADAT
%%%%%%%%%%%%%%%%%%%%%%%%%%
\section*{1. Feladat: Határozza meg a spirálrugó $k^*$ direkciós nyomatékát!}

Vízmérték segítségével az asztallapot vízszintesbe állítottuk.
Asztal sorszáma: 4

$r_{korong}= 6,5 cm$ \\
$m_{edény}=4,6 g$ \\
$\phi_{kezdeti}= 0^\circ$ \\
$m_{golyó}=4,07 g$

\begin{table}[H]
    \centering
    \begin{tabular}{|c|c|c|}        \hline
        Golyók száma[db] & Szögelfordulás\_1 [°] & Szögelfordulás\_2 [°] \\ \hline
         0 & 0 & \\ \hline 
         1 & 42 & 45\\ \hline
         2 & 65 & 60\\ \hline 
         3 & 89 & 82\\ \hline
         4 & 110 & 100\\ \hline 
         5 & 139 & 122\\ \hline
         6 & 157 & 141\\ \hline 
         7 & 185 & 157\\ \hline
         8 & 201 & 183\\ \hline
    \end{tabular}
    \caption{Szögelfordulás mértéke a csapágygolyók számától függően}
    \label{tab:szogelfordulas}
\end{table}

\begin{figure}[H]
    \centering
    \includegraphics[width=0.8\textwidth]{../images/1_feladat.png}
    \caption{Szögkitérés mértékének függése a golyók számától}
    \label{fig:1_feladat_diagram}
\end{figure}

Ellenőrző számításként a direkciós erőt kiszámoltuk: $k^*=1,31 \cdot 10^{-4} \frac{kg m^2}{s^2}$

%%%%%%%%%%%%%%%%%%%%%%%%%%
% 2. FELADAT
%%%%%%%%%%%%%%%%%%%%%%%%%%
\section*{2. Feladat: Határozza meg a rendszer $\beta = \frac{\gamma}{2}$ csillapítási tényezőjét!}

\begin{table}[H]
    \centering
    \begin{tabular}{|c|c|c|}
        \hline
        Szögkitérítés [°] & Idő [s] & Periódusidő [s] \\ \hline
        90 & 12,81 & 1,28\\ \hline 
        100 & 13,06 & 1,30\\ \hline
        110 & 12,85 & 1,29\\ \hline 
        120 & 12,84 & 1,28\\ \hline
        130 & 12,83 & 1,28\\ \hline 
        140 & 12,81 & 1,28\\ \hline
        150 & 12,79&  1,28\\ \hline 
         &  & \\ \hline
    \end{tabular}
    \caption{Periódusidő a kezdeti szögkitérítéstől függően}
    \label{tab:periodusido_2}
\end{table}

20 lengés után nézzük a maximális amplitúdót.
$\phi_0 = 90^\circ$
%videó!!!!!

Átlagolásból: $T = 1,28 \pm 0,01 \text{s}$

\begin{table}[H]
    \centering
    \begin{tabular}{|c|c|}
        \hline
        Mérések száma [db] & Maximális amplitúdó [°] \\ \hline
         1 & 17 \\ \hline 
         2 & 15 \\ \hline
         3 & 17 \\ \hline 
         4 & \\ \hline
         5 & \\ \hline 
         6 & \\ \hline
         & \\ \hline 
         & \\ \hline
    \end{tabular}
    \caption{Maximális amplitúdó}
    \label{tab:max_amplitudo}
\end{table}

Átlagosan: 15,6 $\pm$ 1,1 ° \\
Becsült csillapítási tényező: $3,4 \cdot 10^{-3} s^{-1}$

%%%%%%%%%%%%%%%%%%%%%%%%%%
% 3. FELADAT
%%%%%%%%%%%%%%%%%%%%%%%%%%
\section*{3. Feladat: Határozza meg a torziós asztal tehetetlenségi nyomatékát három különböző módon}

Mérőpár mindkét tagja mérje meg a paramétereket!

\renewcommand{\thesubsection}{\alph{subsection})} % Changes 1.1 to a)
\subsection{Számítással az asztal tömegének és sugarának ismeretében}

Mérőpár mindkét tagja mérje meg a paramétereket!

$\rho_{alumínium} = 2700 \frac{kg}{m^3}$

\begin{table}[H]
    \centering
    \renewcommand{\arraystretch}{1.2}
    \begin{tabular}{|l|c|c|c|}
        \hline
        Mérőpár & $R_{tárcs}$ [m] & $d_{tárcsa}$ & $\Theta$ [kg m$^2$] \\ \hline
        Luca& 0,13 & 0,004 & $2,23 \cdot 10^{-3}$\\ \hline
        Csongor& 0,13 & 0,004 & $2,23 \cdot 10^{-3}$\\ \hline
    \end{tabular}
    \caption{Asztal tehetetlenségi nyomatékának meghatározása}
    \label{tab:hotan_2}
\end{table}


\subsection{A rugó direkciós nyomatékának, a lengésidőnek és a csillapítási tényezőnek az ismeretében}

\[
\theta = \frac{k^*}{4\pi^2} \cdot T^2
\]

$\theta= 1,93 \cdot 10^{-5} \text{kg m}^2$

\subsection{Egy ismert tehetetlenségi nyomatékú tárcsa felhasználásával!}

$m_{korong} = 51,42 \text{g}$ \\
$r_{korong} = 2,5 \text{cm}$ \\
$\theta_{korong} = 1,6 \cdot 10^{-5} \text{kg m}^2$ \\
$m_{rézcsap} = 2,2 \text{g}$

\begin{table}[H]
    \centering
    \begin{tabular}{|c|c|c|}
        \hline
        Szögkitérítés [°] & Idő [s] & Periódusidő [s] \\ \hline
        90 & 13,02 & 1,30\\ \hline 
        100 & 13,07 & 1,31\\ \hline
        110 & 13,08 & 1,31\\ \hline 
        120 & 13,11 & 1,31\\ \hline
        130 & 12,96 & 1,30\\ \hline 
        140 & 12,97 & 1,30\\ \hline
        150 & 13,17 & 1,32\\ \hline
    \end{tabular}
    \caption{Periódusidő a kezdeti szögkitérítéstől függően, ismert tehetetlenséginyomatékú koronggal}
    \label{tab:periodusido_3}
\end{table}

Átlagolásból: $T' = 1,30 \pm 0,01 \text{s}$

Ezekből: $\Theta_{asztal} = 5,1 \cdot 10^{-4} \pm 0,03 \text{kg m}^2$

Most ezt az eredményt úgy számoltuk ki, hogy a rézcsap tömegét nem számoltuk bele amikor kiszámoltuk a korong tehetetlenségi nyomatékát.
Kis eltérést mutat az előző eredményekhez képest.

%%%%%%%%%%%%%%%%%%%%%%%%%%
% 4. FELADAT
%%%%%%%%%%%%%%%%%%%%%%%%%%
\section*{4. Feladat: Igazolja a Steiner-tételt!}

Steiner-tétel: $\theta' = \theta + mr^2.$

$m_{piros csap}= 2,08 \text{kg}$ \\
$r_{korong-középpont}= \text{m}$ \\
$T_{lengési}= \text{s}$

Első mérési pont: \\
1. mérés: $r_{korong-középpont}= \text{m}$, $T_{lengési}= \text{s}$

\begin{figure}[H]
    \centering
    \includegraphics[width=0.8\textwidth]{4_feladat.png}
    \caption{T-r függés illesztése}
    \label{fig:4}
\end{figure}

A direkciós állandóra $6,74 \cdot 10^{-3} \text{Nm}$ adódott.

%%%%%%%%%%%%%%%%%%%%%%%%%%
% 5. FELADAT
%%%%%%%%%%%%%%%%%%%%%%%%%%

\section*{5. Feladat: Határozza meg egy inhomogén tömegeloszlású lemezből készült minta tehetetlenségi nyomatékát a súlypontján átmenő és a lemez síkjára merőleges tengelyre vonatkozóan!}

\[
T'^2_{\max} = \frac{4\pi^2}{k^*} \left[ \theta + \theta_x + m(r_0 + r_1)^2 \right],
\]

\[
T'^2_{\min} = \frac{4\pi^2}{k^*} \left[ \theta + \theta_x + m(r_0 - r_1)^2 \right],
\]

%fotó, minta!!!!

$m_{minta}= 0,387 \text{kg}$
A mintát úgy rögzítjük a 30°-os jelölésnél, hogy a csúcsnál lévő szöfelezővel essen egyvonalba az osztás.

$r_{1}= 5 \text{cm}$
Azért ezt a lyukat választottuk, mert ezt van a középponttól legtávolabb.

Hegyeszögnél lévő lyuknál rögzítettük:

\begin{table}[H]
    \centering
    \begin{tabular}{|c|c|}
        \hline
        Szögelfordulás [°] & 10 lengés lengési ideje [s] \\ \hline
         0 & 17,91 \\ \hline 
         30 & 17,71 \\ \hline
         60 & 16,99 \\ \hline 
         90 & 15,93 \\ \hline
         120 & 14,75 \\ \hline 
         150 & 14,03 \\ \hline
         180 & 13,66 \\ \hline
         210 & 13,99 \\ \hline
         240 & 14,60 \\ \hline
         270 & 15,88 \\ \hline
         300 & 16,77 \\ \hline
         330 & 17,58\\ \hline
        
    \end{tabular}
    \caption{Lengési idők az első lyukkal}
    \label{tab:lengesiido_1.1}
\end{table}

szögnél lévő lyuknál rögzítettük:

\begin{table}[H]
    \centering
    \begin{tabular}{|c|c|}
        \hline
        Szögelfordulás [°] & 10 lengés lengési ideje [s] \\ \hline
         0 & 16,63 \\ \hline 
         30 & 17,28 \\ \hline
         60 & 14,91\\ \hline 
         90 & 17,25 \\ \hline
         120 & 16,60 \\ \hline 
         150 & 15,59 \\ \hline
         180 & 14,59 \\ \hline
         210 & 14,06 \\ \hline
         240 & 13,99 \\ \hline
         270 & 14,18 \\ \hline
         300 & 15,07 \\ \hline
         330 & 15,93 \\ \hline
        
    \end{tabular}
    \caption{Lengési idők a második lyukkal}
    \label{tab:lengesiido_1.2}
\end{table}

$r_{2}= 2 \text{cm}$
Azért ezta szöget választottuk, mert ez van a középponthoz legközelebb

Hegyesszögnél lévő furatnál rögzítve:

\begin{table}[H]
    \centering
    \begin{tabular}{|c|c|}
        \hline
        Szögelfordulás [°] & 10 lengés lengési ideje [s] \\ \hline
         0 & 15,09 \\ \hline 
         30 & 15,13 \\ \hline
         60 & 14,59\\ \hline 
         90 & 14,45\\ \hline
         120 & 13,84 \\ \hline 
         150 & 13,48 \\ \hline
         180 & 13,44 \\ \hline
         210 & 13,41 \\ \hline
         240 & 13,81 \\ \hline
         270 & 14,25 \\ \hline
         300 & 14,73 \\ \hline
         330 & 14,98 \\ \hline
        
    \end{tabular}
    \caption{Lengési idők a második lyukkal}
    \label{tab:lengesiido_2.1}
\end{table}

Tompaszögnél lévő luyknál rögzítve:

\begin{table}[H]
    \centering
    \begin{tabular}{|c|c|}
        \hline
        Szögelfordulás [°] & 10 lengés lengési ideje [s] \\ \hline
         0 & 14,47 \\ \hline 
         30 & 14,84 \\ \hline
         60 & 14,92 \\ \hline 
         90 & 14,72 \\ \hline
         120 & 14,33 \\ \hline 
         150 & 14,05 \\ \hline
         180 & 13,75 \\ \hline
         210 & 13,40 \\ \hline
         240 & 13,27 \\ \hline
         270 & 13,45 \\ \hline
         300 & 13,71 \\ \hline
         330 & 14,11\\ \hline
        
    \end{tabular}
    \caption{Lengési idők a második lyukkal}
    \label{tab:lengesiido_2.2}
\end{table}

$r_{3}= 3 \text{cm}$

Hegyesszögnél lévő furatnál rögzítve:

\begin{table}[H]
    \centering
    \begin{tabular}{|c|c|}
        \hline
        Szögelfordulás [°] & 10 lengés lengési ideje [s] \\ \hline
         0 & 15,89 \\ \hline 
         30 & 15,62 \\ \hline
         60 & 15,41 \\ \hline 
         90 &  \\ \hline
         120 &  \\ \hline 
         150 &  \\ \hline
         180 &  \\ \hline
         210 &  \\ \hline
         240 &  \\ \hline
         270 &  \\ \hline
         300 &  \\ \hline
         330 &  \\ \hline
        
    \end{tabular}
    \caption{Lengési idők a harmadik lyukkal}
    \label{tab:lengesiido_3.1}
\end{table}

Tompaszögnél lévő lyuknál rögzítve:

\begin{table}[H]
    \centering
    \begin{tabular}{|c|c|}
        \hline
        Szögelfordulás [°] & 10 lengés lengési ideje [s] \\ \hline
         0 &  \\ \hline 
         30 &  \\ \hline
         60 &  \\ \hline 
         90 &  \\ \hline
         120 &  \\ \hline 
         150 &  \\ \hline
         180 &  \\ \hline
         210 &  \\ \hline
         240 &  \\ \hline
         270 &  \\ \hline
         300 &  \\ \hline
         330 &  \\ \hline
        
    \end{tabular}
    \caption{Lengési idők a harmadik lyukkal}
    \label{tab:lengesiido_3.2}
\end{table}

A tehetetlenségi noymaték kiszámítására nem maradt időnk. Az időhiány miatt amúgy se maradt időnk az összes lehetséges rögzítési pontnál megmérni a lengési időket, de a középponthoz legközelebbi, és attól legtávolabbi pontban megmértük. Egy harmadik pontotban is elkezdtünk mérni, de nem maradt időnk befejezni.



\end{document}
